\chapter{orthogonal vectors and matrices}

\begin{definition}{adjoint(hermitian conjugate)}
\[
    ^*:
    A = \begin{bmatrix} a_{11} & a_{12} \\ a_{21} & a_{22}   \\a_{31} & a_{32}  \end{bmatrix}
    \mapsto A^* =  \begin{bmatrix} \overline{a_{11}} & \overline{a_{21}}  &\overline{a_{31}} \\ \overline{a_{12}} & \overline{a_{22}} & \overline{a_{32}}   \end{bmatrix}
    \]
(where each overline means complex conjuate.)
\end{definition}
\begin{remark}
    Can easily comfirm:
    $$
    (AB)^* = B^* A^*
    $$
\end{remark}



\begin{definition}{standard inner product and norm on $\bC^m$}
The \textbf{standard inner product}:
$$
<x,y> := x^* y  = \sum_{i=1}^m \overline{x_i} y_i
$$
The \textbf{standard norm}:
$$
||x|| : =  \sqrt{x^* x}
$$
\end{definition}
\begin{remark}
$$
\sqrt{x^* x} = \sqrt{\sum_{i=1}^m |x_i|^2}
$$, 其中 $|x_i| = |a+ bi| = \sqrt{a^2 + b^2}$, \\
\textbf{Note: It is actually $\sqrt{\overline{x_i} x_i}$, but not $\sqrt{x_i^2}$, 这是因为 $\sqrt{\overline{x_i} x_i}$ 等于这个 complex scalar isomorphic 到 $\bR^2$ 上的 Euclidean norm,}\\ 
这是因为两个复数相乘等于长度相乘幅角相加, 而 conjugate 的幅角是相反的, \textbf{所以 conjugate 之间相乘等于幅角相互抵消, 结果在 positive real axis 上, 取开方得到长度.} \\
而平方得到的则是一个转两次的幅角.
\end{remark}
\begin{remark}
    Can easily verifies:
    $$
    <a,b>_{\bC} = \overline{<b,a>_{\bC}}
    $$
    且注意: By bilinearity, \textbf{对于 $a$ 上的 scaling 在拿出 inner product 外后要进行 conjugate.}
\end{remark}

\begin{definition}{orthogonal, orthonomal vectors}
    Say $x, y \in \bC^m$ 是 orthogonal vectors, if $x^* y = 0$.\\
    Say $S \sub \bC^m$ 是 orthogonal 的, 如果其中的 vectors 相互 orthogonal.\\
    Say $S \sub \bC^m$ 是 orthonomal 的, 如果其中的 vectors 相互 orthogonal, 并且每个 vector 的 norm 都是 1.
\end{definition}

\begin{theorem}{orthogonal $\implies$ lin.ind}
    orthogonal 的 vectors 一定 linearly independent.
\end{theorem}
\begin{proof}
    trivial.
\end{proof}
\begin{corollary}
    orthogonal 的 $\dim(V)$ 个 vectors 一定是 $V$ 的一个 basis.
\end{corollary}


\section{decomposing vector by an orthonormal set}
\begin{theorem}{decomposing vector by an orthonormal set}
给定 $\bC^m$ 中的一个 orthonormal set $\{q_1, \cdots, q_n\}$ (by inner product $<\cdot>$, 这里以 standard complex inner product 为例), 
\\对于一个 arbitrary vector $v$, 我们 define:
$$
r:= v - \sum_{i=1}^n <q_i, v>   q_i
$$
Claim: 这个 \textbf{$r$ is orthogonal to $\{q_1, \cdots, q_n\}$,} 即我们把这个 $v$ 分解成了在这个 orthonormal set 上的投影与一个和它们都正交的 vector.
\end{theorem}
\begin{proof}
    注意: 由于 $q_i$ 都是 unit vectors, $$<q_i, v> q_i = \frac{v}{||v||} \cos \alpha = proj_{q_i}(v)$$\textbf{ is the projection of $v$ onto the direction of $q_i$.}\\
我们在两边取和 $q_i$ 的 inner product, for each $i$. 由 linearity 可拆开, 由 orthgonality 可得到:
$$
<q_i,r> = <q_i,v> - <q_i,v><q_i,q_i>
$$
并且由于 $q_i$ 是 unit vector, 得到 $<q_i,q_i>=1$, 从而右边为 0.
\end{proof}
\begin{remark}
    如果 $n=m$, 那么 $r=0$, 我们把 arbitrary vector 分解成了 $\{q_1, \cdots, q_n\}$ 方向上的向量, 相当于对它进行了 change of basis. 这个 change of basis matrix 就等于 $\begin{bmatrix} q_1  & \cdots  & q_n \end{bmatrix}^T$.
\end{remark}


 对于 unit vector $w$, 我们刚才已经展示了一个 arbitrary vector $v$ 在它上面的 projection 是:$$proj_w (v) = <w, v> w  $$
 现在我们引入另一个形式的 projection 表达: projection matrix

\begin{theorem}{projection matrix}
对于任意的 \textbf{unit vector $w$}, we have
$$proj_w (v) = (w \otimes w^*) v $$
其中 $w\otimes w^*$ is called the \textbf{projection matrix} onto $w$.
\end{theorem}
\begin{proof}
    In md.
\textbf{    Notice that this matrix is rank 1.}
\end{proof}


\begin{definition}{unitrary matrix}
一个 square matrix $Q \in \bC^{m\times m}$ 被称为 unitrary 的, if $Q^* = Q^{-1}$.
\end{definition}

\begin{remark}
    unitrary: 即 $QQ^* = Q^*Q = I$. \\
    \textbf{In real case, 它被称为 orthogonal matrix.}
\end{remark}


\begin{theorem}{unitrary matrix 的充要条件}
$Q \in \bC^{m\times m}$ is unitrary $\Longleftrightarrow $ \textbf{its columns are orthonormal} $\Longleftrightarrow$ \textbf{its rows are orthonormal} 
\end{theorem}
\begin{proof}
    显然, 因为 unitrary $\Longleftrightarrow QQ^* = Q^*Q = I$ $\Longleftrightarrow  <q_i, q_j> = \delta_{ij}$
\end{proof}



\begin{theorem}{unitrary transfromation preserves inner product and length}
如果 $Q \in \bC^{m\times m}$ is unitrary, 那么对于任意的 $x,y \in \bC^m$, 都有:
$$
(Qx)^* (Qy) = x^*y
$$
并且自然得到 $$||Qx|| = ||x||
$$
\end{theorem}
\begin{proof}
Follows from: $(AB)^*= B^* A^*$:  $$(Qx)^* (Qy) = x^*Q^*Qy= x^*y
$$
\end{proof}
\begin{remark}
    并不 preserve 自定义的 inner product, 只 \textbf{preserve standard inner product 和 standard norm.}
\end{remark}



